\section{引言}

\subsection{背景与动机}

近年来，边缘AI应用快速增长。用户希望在本地设备上运行大语言模型（LLM）和图像生成模型，以保护隐私、降低延迟。典型场景包括：边与LLM对话边生成图像，或在本地完成文档理解与配图创作。然而，边缘设备资源严重受限——主流笔记本仅配备8GB显存，而7B参数的LLM（Q4\_K\_M量化后模型权重约4.7GB，加上KV Cache约1-2GB，总计约6-7GB）与SDXL Turbo图像模型（fp16精度约6.7GB）无法同时驻留显存。这种资源约束与多模态需求的矛盾，迫切需要轻量级的任务调度系统。

\subsection{核心挑战}

\textbf{显存硬约束}。8GB显存无法同时容纳LLM（模型权重约4.7GB + KV Cache约1-2GB，总计约6-7GB）和图像模型（约6.7GB），传统GPU隔离方案不可行。当图像生成任务到达时，必须立即释放LLM占用的显存，否则触发CUDA OOM。

\textbf{任务异构性}。不同任务具有截然不同的特性，如表~\ref{tab:task_characteristics}所示：

\begin{table}[htbp]
\caption{任务特性对比}
\label{tab:task_characteristics}
\centering
\begin{tabular}{@{}lcccc@{}}
\toprule
任务类型 & 资源需求 & 延迟敏感度 & 可排队性 & 运行语义 \\
\midrule
LLM对话 & GPU/CPU & 高（流式） & 低 & 常驻交互 \\
图像生成 & GPU必需 & 中（端到端） & 高 & 一次性任务 \\
\bottomrule
\end{tabular}
\end{table}

LLM需要持续低延迟响应，而图像生成可以排队等待。这种异构性要求调度系统支持非对称的资源迁移策略。

\subsection{现有方案的局限}

现有方案主要面向两个方向：（1）数据中心分布式调度（如Ray、Kubernetes），启动时间在秒级到分钟级，内存占用GB级别，不适合资源受限的边缘设备；（2）单任务推理优化（如vLLM、FlexGen），假设有充足的GPU资源，缺乏在显存不足时迁移到CPU的能力。

边缘设备上的多模态AI调度问题尚未得到充分研究。现有方案无法解决8GB显存设备上LLM与图像生成模型共存的资源竞争问题，缺乏轻量级、低开销的调度解决方案。EdgeSched填补了这一空白，提供了专为边缘设备设计的多模态任务调度机制。

\subsection{我们的贡献}

本文提出EdgeSched，一个面向边缘设备的多模态AI调度系统，核心贡献包括：

\begin{enumerate}
\item \textbf{轻量级架构}。调度器启动时间毫秒级，内存占用约27MB，适合资源受限的边缘设备。

\item \textbf{非对称动态资源迁移}。针对显存硬约束与用户体验软约束的非对称性，设计紧急模式（GPU→CPU立即释放）与原子热替换模式（CPU→GPU预加载），将服务中断时间从秒级压缩到毫秒级。

\item \textbf{完整的任务生命周期管理}。支持优先级队列、背压控制、取消语义和崩溃恢复（<12ms），提供可预测的调度行为。

\item \textbf{开源实现与验证}。在RTX 5070 Laptop（8GB显存）上验证，混合负载下成功率达100\%，服务中断时间<300ms。
\end{enumerate}

EdgeSched填补了边缘设备上多模态AI调度系统的空白，为资源受限场景下的LLM与图像生成共存提供了实用解决方案。
